\documentclass{beamer}

\usepackage[T1]{fontenc}
\usepackage[utf8]{inputenc}
\usepackage{lmodern}
\usepackage{amsmath,amssymb,amsfonts}
\usepackage{physics}
\usepackage{bm}
\usepackage{mathtools}
\usepackage{quantikz}

\title{Quantum Fourier Transform and Quantum Phase Estimation}
\author{Morten Hjorth-Jensen}
\date{March 25, 2026}

\begin{document}

%------------------------------------------------
\frame{\titlepage}

%------------------------------------------------
\begin{frame}{Outline}
\begin{enumerate}
\item Review of the Quantum Fourier Transform (QFT)
\item Mathematical derivation of the QFT circuit
\item Quantum Phase Estimation (QPE): goal and idea
\item Step-by-step derivation of QPE
\item Bit-string interpretation and precision
\item Connection to Hamiltonian simulation
\item Connection to Variational Quantum Eigensolvers (VQE)
\item Summary and outlook
\end{enumerate}
\end{frame}

%================================================
\begin{frame}{The Quantum Fourier Transform}
For \(N=2^n\), the Quantum Fourier Transform is the unitary map
\[
\mathrm{QFT}_N \ket{x}
=
\frac{1}{\sqrt N}
\sum_{k=0}^{N-1}
e^{2\pi i xk/N}\ket{k},
\qquad x=0,\dots,N-1.
\]

This is the quantum analogue of the discrete Fourier transform.

The QFT acts on amplitudes, not on a classical list of numbers.
\end{frame}

%================================================
\begin{frame}{Binary Representation}
Write
\[
x = x_1 x_2 \dots x_n
\]
in binary, with
\[
x = \sum_{j=1}^n x_j 2^{n-j},
\qquad x_j\in\{0,1\}.
\]

We also define the binary fraction
\[
0.x_j x_{j+1}\dots x_n
=
\frac{x_j}{2}+\frac{x_{j+1}}{2^2}+\cdots+\frac{x_n}{2^{n-j+1}}.
\]

This notation is essential for factorizing the QFT output.
\end{frame}

%================================================
\begin{frame}{Product Form of the QFT}
A standard derivation gives
\[
\mathrm{QFT}\ket{x_1x_2\dots x_n}
=
\frac{1}{2^{n/2}}
\left(\ket0+e^{2\pi i\,0.x_n}\ket1\right)
\otimes
\left(\ket0+e^{2\pi i\,0.x_{n-1}x_n}\ket1\right)
\otimes\cdots
\]
\[
\cdots\otimes
\left(\ket0+e^{2\pi i\,0.x_1x_2\dots x_n}\ket1\right).
\]

This shows that the QFT output factorizes into one-qubit states with progressively longer binary-fraction phases.
\end{frame}

%================================================
\begin{frame}{Elementary Gates in the QFT}
The QFT circuit is built from

\textbf{Hadamard gates}
\[
H = \frac{1}{\sqrt 2}
\begin{pmatrix}
1 & 1\\
1 & -1
\end{pmatrix}
\]

and controlled phase rotations
\[
R_k =
\begin{pmatrix}
1 & 0\\
0 & e^{2\pi i/2^k}
\end{pmatrix}.
\]

These gates generate the binary-fraction phases appearing in the product form.
\end{frame}

%================================================
\begin{frame}{Two-Qubit QFT Circuit}
\centering
\begin{quantikz}
\lstick{$q_0$} & \gate{H} & \ctrl{1} & \swap{1}   & \qw \\
\lstick{$q_1$} & \qw      & \gate{R_2} & \targX{}{} & \qw
\end{quantikz}

\vspace{0.5em}

The final SWAP corrects the reversed qubit order in the product form.
\end{frame}

%================================================
\begin{frame}{Three-Qubit QFT Circuit}
\centering
\begin{quantikz}
\lstick{$q_0$} & \gate{H} & \ctrl{1}   & \ctrl{2}   & \qw      & \qw        & \swap{2}   & \qw \\
\lstick{$q_1$} & \qw      & \gate{R_2} & \qw        & \gate{H} & \ctrl{1}   & \qw        & \qw \\
\lstick{$q_2$} & \qw      & \qw        & \gate{R_3} & \qw      & \gate{R_2} & \targX{}{} & \qw
\end{quantikz}

\vspace{0.7em}

This shows the characteristic pattern: Hadamards, controlled phases, then bit reversal.
\end{frame}

%================================================
\begin{frame}{Inverse QFT}
The inverse QFT is obtained by reversing the circuit and replacing each
\[
R_k \mapsto R_k^\dagger.
\]

Thus the inverse QFT is also efficient.

This inverse transform is the form that appears directly in Quantum Phase Estimation.
\end{frame}

%================================================
\begin{frame}{Quantum Phase Estimation: Problem Statement}
Suppose \(U\) is a unitary operator and \(\ket{\psi}\) is one of its eigenstates:
\[
U\ket{\psi}=e^{2\pi i \phi}\ket{\psi},
\qquad 0\le \phi <1.
\]

The goal of Quantum Phase Estimation (QPE) is to estimate the phase \(\phi\).

Equivalently, QPE extracts the eigenvalue of \(U\) in phase form.
\end{frame}

%================================================
\begin{frame}{Why QPE Matters}
QPE is central because many physical problems can be mapped to eigenvalue estimation.

Examples:
\begin{itemize}
\item energy estimation in quantum chemistry,
\item eigenvalue problems in many-body physics,
\item Hamiltonian simulation,
\item Shor's algorithm and period finding.
\end{itemize}

If we can encode a Hamiltonian into a unitary, then QPE can extract its eigenvalues.
\end{frame}

%================================================
\begin{frame}{Registers Used in QPE}
QPE uses two registers:

\begin{itemize}
\item a control register of \(n\) qubits, initialized in \(\ket{0}^{\otimes n}\),
\item a system register, initialized in the eigenstate \(\ket{\psi}\).
\end{itemize}

So the initial state is
\[
\ket{0}^{\otimes n}\ket{\psi}.
\]

The control register stores the estimated phase bits.
\end{frame}

%================================================
\begin{frame}{Step 1: Create a Uniform Superposition}
Apply Hadamard gates to all control qubits:
\[
\ket{0}^{\otimes n}
\mapsto
\frac{1}{2^{n/2}}
\sum_{k=0}^{2^n-1}\ket{k}.
\]

So the total state becomes
\[
\frac{1}{2^{n/2}}
\sum_{k=0}^{2^n-1}
\ket{k}\ket{\psi}.
\]

This prepares the control register as a coherent superposition of all binary numbers \(k\).
\end{frame}

%================================================
\begin{frame}{Step 2: Controlled Powers of \(U\)}
Now apply controlled powers of \(U\):
\[
U^{2^0}, U^{2^1}, U^{2^2}, \dots, U^{2^{n-1}}.
\]

Each control qubit determines whether the corresponding power is applied.

Because \(\ket{\psi}\) is an eigenstate,
\[
U^{2^j}\ket{\psi} = e^{2\pi i 2^j \phi}\ket{\psi}.
\]

Hence each controlled unitary writes phase information into the control register.
\end{frame}

%================================================
\begin{frame}{Step 2: State After Controlled Powers}
If the control register basis state is \(\ket{k}\), then the controlled operations apply \(U^k\), giving
\[
U^k\ket{\psi} = e^{2\pi i k\phi}\ket{\psi}.
\]

Therefore the combined state becomes
\[
\frac{1}{2^{n/2}}
\sum_{k=0}^{2^n-1}
e^{2\pi i k\phi}\ket{k}\ket{\psi}.
\]

The system register factors out, and the control register now contains the phase information.
\end{frame}

%================================================
\begin{frame}{Key Observation}
The control register is in the state
\[
\frac{1}{2^{n/2}}
\sum_{k=0}^{2^n-1} e^{2\pi i k\phi}\ket{k},
\]
which is exactly a Fourier-type state.

This is why the inverse QFT is the correct next step:
it converts this phase-encoded superposition into a computational basis approximation of \(\phi\).
\end{frame}

%================================================
\begin{frame}{Step 3: Apply the Inverse QFT}
Applying \(\mathrm{QFT}^{-1}\) to the control register gives
\[
\mathrm{QFT}^{-1}
\left(
\frac{1}{2^{n/2}}
\sum_{k=0}^{2^n-1} e^{2\pi i k\phi}\ket{k}
\right).
\]

If
\[
\phi = \frac{m}{2^n}
\]
exactly for some integer \(m\), then this becomes
\[
\ket{m}.
\]

So the control register stores the exact \(n\)-bit binary expansion of \(\phi\).
\end{frame}

%================================================
\begin{frame}{Exact QPE Case}
Suppose
\[
\phi = 0.\phi_1\phi_2\dots\phi_n
=
\frac{m}{2^n}.
\]

Then after inverse QFT, the control register becomes
\[
\ket{\phi_1\phi_2\dots\phi_n}.
\]

Thus measuring the control register returns the exact binary digits of \(\phi\).

This is the cleanest case of QPE.
\end{frame}

%================================================
\begin{frame}{Step-by-Step Derivation with Binary Fractions}
Let
\[
\phi = 0.\phi_1\phi_2\dots\phi_n\dots
\]

Then one can show that the control-register state after controlled powers factorizes as
\[
\frac{1}{2^{n/2}}
\left(\ket0 + e^{2\pi i\, 2^{n-1}\phi}\ket1\right)
\otimes
\left(\ket0 + e^{2\pi i\, 2^{n-2}\phi}\ket1\right)
\otimes \cdots
\]
\[
\cdots \otimes
\left(\ket0 + e^{2\pi i\, \phi}\ket1\right).
\]

Because
\[
2^{n-1}\phi = \phi_1.\phi_2\phi_3\dots,
\]
only the fractional parts matter in the phases, so each qubit encodes one successive binary suffix.
\end{frame}

%================================================
\begin{frame}{How the Bits Emerge}
Using the binary expansion of \(\phi\), the control-register state becomes
\[
\frac{1}{2^{n/2}}
\left(\ket0 + e^{2\pi i\,0.\phi_n}\ket1\right)
\otimes
\left(\ket0 + e^{2\pi i\,0.\phi_{n-1}\phi_n}\ket1\right)
\otimes \cdots
\]
\[
\cdots \otimes
\left(\ket0 + e^{2\pi i\,0.\phi_1\phi_2\dots\phi_n}\ket1\right),
\]
which has exactly the structure of a QFT state.

Therefore applying \(\mathrm{QFT}^{-1}\) yields the bit string
\[
\phi_1\phi_2\dots\phi_n.
\]
\end{frame}

%================================================
\begin{frame}{QPE Circuit}
\centering
\begin{quantikz}
\lstick{$\ket{0}$}    & \gate{H} & \ctrl{3} & \qw      & \qw      & \gate[3]{\mathrm{QFT}^\dagger} & \meter{} \\
\lstick{$\ket{0}$}    & \gate{H} & \qw      & \ctrl{2} & \qw      & \ghost{\mathrm{QFT}^\dagger}   & \meter{} \\
\lstick{$\ket{0}$}    & \gate{H} & \qw      & \qw      & \ctrl{1} & \ghost{\mathrm{QFT}^\dagger}   & \meter{} \\
\lstick{$\ket{\psi}$} & \qw      & \gate{U^4} & \gate{U^2} & \gate{U} & \qw                          & \qw
\end{quantikz}

\vspace{0.5em}

This example shows a 3-qubit control register and one system register.
\end{frame}

%================================================
\begin{frame}{Worked Example: \(\phi=5/8\)}
Suppose
\[
\phi = \frac{5}{8} = 0.101.
\]

Then the ideal QPE output on a 3-qubit control register is
\[
\ket{101}.
\]

Indeed,
\[
e^{2\pi i \phi} = e^{2\pi i (5/8)}.
\]

After the controlled powers and inverse QFT, the control register collapses to the bit string \(101\) with probability 1.
\end{frame}

%================================================
\begin{frame}{When \(\phi\) Is Not Exactly Representable}
If \(\phi\) is not exactly of the form \(m/2^n\), then QPE does not return one exact bit string.

Instead, the output is a probability distribution peaked near the best \(n\)-bit approximation:
\[
\phi \approx \frac{m}{2^n}.
\]

The measurement outcome gives the nearest binary approximation with high probability.
\end{frame}

%================================================
\begin{frame}{Precision of QPE}
With \(n\) control qubits, QPE resolves phases to approximately
\[
\Delta \phi \sim \frac{1}{2^n}.
\]

Thus the precision improves exponentially with the number of control qubits.

This exponential precision is one of the main reasons QPE is so powerful.
\end{frame}

%================================================
\begin{frame}{From Phases to Energies}
Suppose \(H\) is a Hamiltonian with eigenstate \(\ket{E}\):
\[
H\ket{E} = E\ket{E}.
\]

Define the time-evolution operator
\[
U(t)=e^{-iHt}.
\]

Then
\[
U(t)\ket{E} = e^{-iEt}\ket{E}.
\]

Comparing with
\[
U\ket{\psi}=e^{2\pi i \phi}\ket{\psi},
\]
we identify
\[
\phi = -\frac{Et}{2\pi} \mod 1.
\]

Thus QPE can estimate energies if we can implement \(e^{-iHt}\).
\end{frame}

%================================================
\begin{frame}{Connection to Hamiltonian Simulation}
QPE requires controlled applications of
\[
U^{2^j} = e^{-iH t 2^j}.
\]

So QPE depends directly on the ability to perform Hamiltonian simulation.

In practice this may be done using:
\begin{itemize}
\item Trotter-Suzuki decompositions,
\item linear-combination-of-unitaries methods,
\item qubitization,
\item problem-specific simulation techniques.
\end{itemize}

Thus:
\[
\text{Hamiltonian simulation} + \text{QPE}
\quad \Rightarrow \quad
\text{energy estimation}.
\]
\end{frame}

%================================================
\begin{frame}{Physical Interpretation}
QPE measures the dynamical phase accumulated by an energy eigenstate under time evolution.

The algorithm turns the phase
\[
e^{-iEt}
\]
into a bit string encoding the energy.

So one can view QPE as a highly precise interferometric energy-measurement protocol.
\end{frame}

%================================================
\begin{frame}{Why QPE Is Important in Quantum Chemistry}
In quantum chemistry and many-body physics, the main goal is often to compute eigenenergies of a Hamiltonian.

If we can prepare an approximate eigenstate with sufficient overlap with the true ground state, then QPE can project onto the exact eigenvalue with high precision.

This makes QPE a natural route to:
\begin{itemize}
\item ground-state energies,
\item excited-state energies,
\item spectral properties.
\end{itemize}
\end{frame}

%================================================
\begin{frame}{Connection to VQE}
The Variational Quantum Eigensolver (VQE) and QPE aim at similar physics goals, but in different ways.

\textbf{VQE:}
\begin{itemize}
\item variational,
\item hybrid quantum-classical,
\item optimized by minimizing
\[
E(\bm{\theta}) = \bra{\psi(\bm{\theta})} H \ket{\psi(\bm{\theta})},
\]
\item suitable for noisy near-term devices.
\end{itemize}

\textbf{QPE:}
\begin{itemize}
\item not variational,
\item requires coherent controlled time evolution,
\item provides direct eigenvalue estimation,
\item suited to fault-tolerant quantum computers.
\end{itemize}
\end{frame}

%================================================
\begin{frame}{How VQE and QPE Complement Each Other}
A common long-term strategy is:
\begin{enumerate}
\item Use a variational method such as VQE to prepare a good approximate eigenstate
\[
\ket{\psi_{\mathrm{trial}}}
\]
with large overlap with the true eigenstate.

\item Use that state as the input to QPE.

\item QPE then refines the energy estimate to high precision.
\end{enumerate}

So VQE can be viewed as a state-preparation front end for QPE.
\end{frame}

%================================================
\begin{frame}{Overlap Requirement}
If the initial state is
\[
\ket{\psi_{\mathrm{trial}}} = \sum_j c_j \ket{E_j},
\]
then QPE outputs \(E_j\) with probability
\[
|c_j|^2.
\]

Therefore the success probability depends on the overlap with the desired eigenstate.

This is exactly why improved ans\"atze in VQE are also valuable for QPE-based workflows.
\end{frame}

%================================================
\begin{frame}{QPE vs VQE in Practice}
\textbf{VQE}
\begin{itemize}
\item shallow circuits,
\item many repeated measurements,
\item classical optimization loop,
\item noise tolerant,
\item limited by optimization difficulty and shot cost.
\end{itemize}

\textbf{QPE}
\begin{itemize}
\item deep coherent circuits,
\item requires controlled \(e^{-iHt}\),
\item very precise,
\item ideal for fault-tolerant hardware,
\item often viewed as the long-term gold standard for energy estimation.
\end{itemize}
\end{frame}

%================================================
\begin{frame}{Semiclassical Interpretation}
One may interpret QPE as repeatedly extracting successive binary digits of a phase.

The inverse QFT acts globally, but conceptually it decodes the phase one bit at a time.

This bitwise structure is one reason QPE is so elegant mathematically.
\end{frame}

%================================================
\begin{frame}{Implementation Logic in a Notebook}
A standard notebook implementation of QPE typically proceeds as follows:
\begin{enumerate}
\item choose the number of control qubits,
\item initialize the system register in an eigenstate or approximate eigenstate,
\item apply Hadamards to the control register,
\item apply controlled \(U^{2^j}\),
\item apply inverse QFT,
\item measure the control register,
\item interpret the measurement as an estimate of \(\phi\).
\end{enumerate}

If the notebook also implements the QFT explicitly, it will usually reuse the same Hadamard + controlled-phase + swap pattern.
\end{frame}

%================================================
\begin{frame}{Summary of the Full Picture}
\[
\text{Hamiltonian } H
\quad \longrightarrow \quad
U(t)=e^{-iHt}
\quad \longrightarrow \quad
\text{QPE}
\quad \longrightarrow \quad
E
\]

The QFT is the phase-decoding engine inside QPE.

VQE and QPE are complementary:
\begin{itemize}
\item VQE is a near-term variational energy estimator,
\item QPE is a long-term high-precision eigenvalue estimator,
\item VQE can provide input states for QPE.
\end{itemize}
\end{frame}

%================================================
\begin{frame}{Summary}
In these notes we have seen that:
\begin{itemize}
\item the QFT maps
\[
\ket{x}\mapsto \frac{1}{\sqrt N}\sum_k e^{2\pi i xk/N}\ket{k},
\]
\item QPE uses controlled powers of a unitary and the inverse QFT to extract eigenphases,
\item for Hamiltonian evolution
\[
U(t)=e^{-iHt},
\]
QPE becomes an energy-estimation algorithm,
\item VQE and QPE address related eigenvalue problems with different hardware requirements.
\end{itemize}
\end{frame}

%================================================
\begin{frame}{Possible Next Topics}
Natural follow-ups include:
\begin{itemize}
\item iterative and semiclassical QPE,
\item error and success-probability bounds,
\item Hamiltonian simulation methods,
\item comparison with Bayesian phase estimation,
\item detailed quantum chemistry applications.
\end{itemize}
\end{frame}

\end{document}
